% 第三章 系统需求分析

\chapter{系统需求分析}

\section{应用场景分析}

FreeFlow面向链上音乐资源的发布、访问与消费场景。系统中，创作者负责作品编辑、素材上传、发布配置与链上发行；听众负责资源检索、访问购买、播放、下载与评论；平台运维方负责链环境配置、服务部署与运行维护。因此，系统需要支持从作品发布到资源消费的完整业务闭环，并保证链上状态与链下服务之间的协同配合。

\section{用户角色分析}

系统涉及创作者、听众和平台运维方三类角色，其职责与需求如表\ref{tab:user_roles}所示。

\begin{table}[htbp]
    \linespread{1.5}
    \zihao{5}
    \songti
    \centering
    \caption{用户角色与需求说明}\label{tab:user_roles}
    \begin{tabular}{|l|p{10cm}|}
        \hline
        角色    & 主要需求                                                          \\ \hline
        创作者   & 创建发行草稿、编辑音乐信息、上传音频和封面、生成metadata、设置访问模式、设置收益分账、发布作品到链上、查看发布状态 \\ \hline
        听众    & 搜索链上音乐、查看作品详情、查询访问状态、购买访问权、播放音乐、加入链上音乐库、下载音乐、发表评论             \\ \hline
        平台运维方 & 配置链环境和合约地址、维护SIWE会话与Pinata上传能力、保障资源索引与评论服务可用、处理部署与运行异常        \\ \hline
    \end{tabular}
\end{table}

\section{功能需求分析}

系统需要覆盖的钱包登录、发行、检索、购买、播放与评论等核心功能如表\ref{tab:functional_requirements}所示。

\begin{table}[htbp]
    \linespread{1.5}
    \zihao{5}
    \songti
    \centering
    \caption{系统功能需求表}\label{tab:functional_requirements}
    \begin{tabular}{|l|p{10cm}|}
        \hline
        功能模块        & 需求说明                                           \\ \hline
        钱包登录与用户资料管理 & 支持用户通过钱包签名完成SIWE登录，维护本地用户资料与钱包身份的一致性           \\ \hline
        创作者工作台      & 支持发行草稿创建、编辑、保存、删除和状态展示                         \\ \hline
        素材上传与元数据管理  & 支持音频、封面和元数据上传到Pinata/IPFS，并保存CID和网关地址          \\ \hline
        链上发布与分账     & 支持调用智能合约发布作品，设置访问模式、价格和收益分账地址                  \\ \hline
        链上资源检索      & 支持按标题、歌手、专辑、流派、Token ID、合约地址、交易哈希和IPFS CID检索资源 \\ \hline
        购买与访问验证     & 支持查询作品销售配置、购买访问权、验证ERC-1155访问凭证和记录购买交易         \\ \hline
        播放与链上音乐库    & 支持播放公开或已购买作品，将作品加入链上音乐库，并进行本地缓存和下载             \\ \hline
        评论互动        & 支持登录用户查看和发布作品评论                                \\ \hline
    \end{tabular}
\end{table}

\section{非功能需求分析}

除功能性需求外，系统还需满足安全性、可追溯性、可用性、可维护性和性能等方面要求，具体如表\ref{tab:non_functional_requirements}所示。

\begin{table}[htbp]
    \linespread{1.5}
    \zihao{5}
    \songti
    \centering
    \caption{系统非功能需求表}\label{tab:non_functional_requirements}
    \begin{tabular}{|l|p{10cm}|}
        \hline
        需求类型 & 说明                                                 \\ \hline
        安全性  & SIWE登录需要防止nonce重放，后端受保护接口需要进行会话校验，付费作品访问状态以链上凭证为依据 \\ \hline
        可追溯性 & 作品应通过chainId、合约地址、tokenId、CID和交易哈希进行定位和追踪          \\ \hline
        可用性  & 创作者发布流程应提供明确阶段状态和错误提示，用户购买和播放流程应尽量减少操作步骤           \\ \hline
        可维护性 & 桌面端、后端和合约应按模块划分职责，降低后续扩展成本                         \\ \hline
        性能   & 检索排序应限制候选规模，音频播放应支持Range请求和缓存机制                    \\ \hline
    \end{tabular}
\end{table}

\section{可行性分析}

\subsection{技术可行性}

系统所采用的Electron、React、Node.js、Express、Prisma、Solidity、ethers.js和IPFS/Pinata均具有较成熟的工程实践基础，能够覆盖桌面端开发、链下服务组织、链上交互与资源存储等需求。当前原型已完成主要业务链路，实现方案具备落地条件。

\subsection{经济可行性}

系统主要使用开源技术框架，开发成本集中在实现、测试与部署过程。链上交互使用Sepolia测试网和MetaMask钱包基于SeloliaEth进行,无需主网货币和法币交易。智能合约测试基于Hardhat本地环境,曲目检索方法基于有限规模下完成，整体成本处于可控范围。

\subsection{时间可行性}

当前系统已具备主要功能，后续工作主要包括测试补充、实验整理和论文撰写。在既定周期内完成系统说明与关键验证任务具有现实可行性。
